% Compile with:  pdflatex docs/branch-diff-hybrid-wogrid.tex  (twice)
\documentclass[10pt,aspectratio=169]{beamer}
\usetheme{default}
\usecolortheme{seahorse}
\useoutertheme{infolines}
\setbeamertemplate{navigation symbols}{}
\setbeamerfont{frametitle}{size=\normalsize}
\setbeamertemplate{itemize items}[circle]

\usepackage[T1]{fontenc}
\usepackage{lmodern}
\usepackage{amsmath,amssymb}
\usepackage{booktabs}
\usepackage{array}

% compact bullets
\setlength{\leftmargini}{1.4em}
\setlength{\leftmarginii}{1.2em}
\newcommand{\tit}[1]{\texttt{\small #1}}
\newcommand{\ti}[1]{\texttt{\footnotesize #1}}

\title{\texttt{pochoir}: codebase description}
\subtitle{Branch \texttt{hybrid\_wogrid\_merged} against \texttt{confings\_and\_instructions}}
\author{Reference notes}
\date{\today}

\begin{document}

\begin{frame}[plain]
  \titlepage
\end{frame}

\begin{frame}{Contents}
  \small
  \tableofcontents
\end{frame}

% =====================================================================
\section{Branch relationship}
% =====================================================================

\begin{frame}{Branch topology}
  \begin{itemize}\small
  \item Merge base \ti{f24803c} is the \emph{parent} of \ti{confings\_and\_instructions};
        \ti{hybrid\_wogrid\_merged} is therefore a near-superset.
  \item \ti{hybrid\_wogrid\_merged}: 342 commits, 247 ahead of the other branch.
  \item \ti{confings\_and\_instructions}: 96 commits, 1 commit not in the merged branch
        (\ti{362b3df} ``Add CLAUDE.md''); that file's content is present in both, so the
        file-level diff for \ti{CLAUDE.md} is empty.
  \item Net diff: \textbf{150 files, +27\,572 / $-$324 lines}. No file is deleted; the
        deletions are rewrites inside four pre-existing files.
  \item Head of \ti{hybrid\_wogrid\_merged}: \ti{532e923} ``task13: retarget to 3.7mm pitch,
        9x9, 15cm drift on 0.37/0.0925 grids''.
  \end{itemize}
\end{frame}

\begin{frame}{Diff by file class}
  \small
  \begin{center}
  \begin{tabular}{lr}
    \toprule
    Class & Files \\
    \midrule
    \ti{pochoir/} library modules & 13 \\
    \ti{test/test\_*.py} pytest files & 29 \\
    Shell runners (\ti{scripts/}, \ti{test/}) & 40 \\
    JSON geometry configs & 36 \\
    Jupyter notebooks (\ti{test/}, \ti{toy/}) & 11 \\
    Examination docs (\ti{docs/examinations/wogrid/}) & 6 \\
    Other (\ti{Makefile}, \ti{README\_for\_pixelatedReadout.md}, \ti{.gitignore}, \ti{.beads/}) & 15 \\
    \midrule
    Total & 150 \\
    \bottomrule
  \end{tabular}
  \end{center}
\end{frame}

\begin{frame}[allowframebreaks]{Baseline: \texttt{confings\_and\_instructions}}
  \begin{itemize}\small
  \item The \ti{for\_pix}-lineage pixel branch plus two configs-and-instructions commits
        and \ti{CLAUDE.md}.
  \item Solver: one Laplace Jacobi kernel, \ti{fdm\_generic.stencil()}; \\
        \ti{fdm\_torch.solve(iarr, barr, periodic, prec, epoch, nepochs)}.
  \item No dtype control, no \ti{torch.compile}, no source term, no permittivity,
        no insulator boundary.
  \item CLI: 28 commands (\ti{\_\_main\_\_.py}, 1545 lines). No near/far machinery,
        no field drivers.
  \item Configs carry 9 keys: \ti{HoleRadius}, \ti{PcbWidth}, \ti{GridPotential},
        \ti{CathodePotential}, \ti{pixelPlaneLowEdgePosition}, \ti{pixelPlaneWidth},
        \ti{pixelSize}, \ti{pixelGap}, \ti{Npixels}.
  \item Geometry: circular quarter-holes only; fixed-cell staircase chamfer
        (\ti{trimCorner}, hardcoded 4-cell reach); no FR4, no pad thickness.
  \item Drift: component-wise velocity-field interpolation (\ti{Simple}); no periodic
        transverse wrap; no potential-based drift.
  \end{itemize}
\end{frame}

% =====================================================================
\section{Main scripts: \texttt{scripts/*.sh}}
% =====================================================================

\begin{frame}[allowframebreaks]{\texttt{scripts/} inventory}
  \small
  \begin{center}
  \begin{tabular}{p{0.40\linewidth}p{0.52\linewidth}}
    \toprule
    File & Role \\
    \midrule
    \ti{run-pixel-field.sh} & main runner; hybrid or single-spacing, task13 geometry \\
    \ti{run-for-larpix-v2a-wogrid.sh} & LArPix v2a, no shield grid, single full-depth solve \\
    \ti{run-for-largepix-wgrid.sh} & large pixel \emph{with} circular-hole shield grid \\
    \ti{helpers.sh} & \ti{want} / \ti{want\_file} resume guards \\
    \ti{run-task10b-drift-2cm-gap09-chamf07.sh} & validated 2\,cm drift reference run \\
    \ti{run-task12d-gridsq-30cm.sh} & square-hole shield grid, 30\,cm drift \\
    \ti{run-pixel-field.md}, \ti{run-pixel-field-geometry.md} & command-by-command walkthroughs \\
    \ti{example\_gen\_*.json} (12) & geometry configs read by the above \\
    \bottomrule
  \end{tabular}
  \end{center}
\end{frame}

\begin{frame}{\texttt{helpers.sh} and the resume guard}
  \begin{itemize}\small
  \item \ti{want TARGET CMD...}: if \ti{\$POCHOIR\_STORE/TARGET.npz} exists, print
        \ti{have TARGET} and return; else run \ti{CMD}, verify the file appeared, print
        \ti{made TARGET}, or exit non-zero.
  \item \ti{want\_file TARGET CMD...}: same, for a plain path rather than a store key.
  \item \ti{run-pixel-field.sh} and the two \ti{run-for-*} scripts \emph{override}
        \ti{want} with a multi-key variant: first argument is a space-separated list of
        store keys; the step is skipped only if all exist.
  \item Effect: an interrupted run resumes cleanly against the same store; a step that
        silently produces nothing is a hard error, not a skip.
  \item \ti{POCHOIR\_STORE} is exported by each script (default per-script store folder
        name, overridable as \ti{\$1}); \ti{POCHOIR\_LOG} is re-pointed per phase
        (\ti{pochoir\_driftfield.log}, \ti{pochoir\_weightingfield.log}).
  \end{itemize}
\end{frame}

\begin{frame}[allowframebreaks]{\texttt{run-pixel-field.sh} --- structure}
  \begin{itemize}\small
  \item Usage: \ti{./run-pixel-field.sh [--hybrid yes|no] [STORE\_DIR]};
        anything but \ti{yes}/\ti{no} exits 1. \ti{set -e}; \ti{cd "\$(dirname \$0)"} so
        configs resolve as bare filenames.
  \item \textbf{PART A --- drift field}: \ti{pochoir field-solve --field drift}
        (\ti{--hybrid yes} $\rightarrow$ coarse+near+fine stitch; \ti{--hybrid no}
        $\rightarrow$ one full-depth solve). Writes \ti{potential/drift3d}.
  \item \textbf{PART B --- velocity and paths}, identical in both modes:
    \begin{itemize}
    \item \ti{pochoir velo --temperature '87.0*K'} $\rightarrow$ \ti{velocity/drift3d};
    \item \ti{pochoir starts -m yes -c \$dcfg\_fine --plot} $\rightarrow$ \ti{starts/drift3d};
    \item \ti{pochoir drift --interp-order linear ... '0*us,200*us,0.05*us'}
          $\rightarrow$ \ti{paths/drift3d}.
    \end{itemize}
  \item \textbf{PART C --- weighting field and current}:
        \ti{field-solve --field weighting} $\rightarrow$ \ti{potential/weight3d};
        \ti{induce-pixel --npixels 2 --config \$wcfg\_fine --plot} $\rightarrow$
        \ti{current/induced\_current}.
  \item \ti{date} is printed after the drift solve, the weighting solve and the induce
        step; run ends with \ti{=== DONE ===}.
  \end{itemize}
\end{frame}

\begin{frame}{\texttt{run-pixel-field.sh} --- SIZES block}
  \begin{itemize}\small
  \item Four configs, bare filenames: \ti{example\_gen\_pcb\_drift\_pixel\_task13\_\{coarse,fine\}.json},
        \ti{example\_gen\_pixel\_with\_grid\_task13\_\{coarse,fine\}.json}.
  \item Shapes are an explicit hand-maintained table; \ti{field-solve} is invoked with
        \ti{--domain no}, so the table is the authority and no shape is re-derived.
  \item \ti{--domain no} also bypasses \ti{hybrid\_iterate.\_cells()}, the only check that
        would refuse a geometry the spacings cannot represent.
  \item Scalars: \ti{interface='39.96*mm'}, \ti{coarse\_spacing=0.37},
        \ti{fine\_spacing=0.0925}, \ti{spacing=0.0925}, \ti{precision=2e-8}.
  \item The \ti{single} row and the \ti{fine} row are the same grid by design, kept as two
        labelled variables rather than one.
  \end{itemize}
\end{frame}

\begin{frame}[allowframebreaks]{\texttt{run-pixel-field.sh} --- shape table (head commit)}
  \small
  Geometry: pitch 3.7\,mm (\ti{pixelSize} 2.2 + \ti{pixelGap} 1.5), \ti{Npixels} 9,
  \ti{driftZDepth} 159.9, interface 39.96\,mm.
  \begin{center}
  \begin{tabular}{llll}
    \toprule
    Grid & Spacing / extent & drift & weighting \\
    \midrule
    coarse & 0.37\,mm, full depth      & \ti{10,10,434}  & \ti{90,90,434} \\
    near   & 0.0925\,mm, to interface  & \ti{40,40,433}  & \ti{360,360,433} \\
    fine   & 0.0925\,mm, full depth    & \ti{40,40,1733} & \ti{360,360,1733} \\
    single & 0.0925\,mm (\ti{--hybrid no}) & \ti{40,40,1733} & \ti{360,360,1733} \\
    \bottomrule
  \end{tabular}
  \end{center}
  \begin{itemize}\small
  \item depth 160.21\,mm $=\lceil 159.9/0.37\rceil = 433$ coarse cells $=433\times4=1732$
        fine cells exactly.
  \item near 0--39.96\,mm $=432$ fine cells; 39.96\,mm $=108$ coarse cells, so the seam is
        a plane of \emph{both} grids.
  \item pitch 3.7\,mm $\rightarrow$ 10 coarse / 40 fine cells;
        probe $9\times3.7=33.3$\,mm $\rightarrow$ 90 coarse / 360 fine cells.
  \item Weighting fine grid $\approx225$\,M nodes (was $\approx77$\,M at \ti{220,220,1601}).
  \end{itemize}
\end{frame}

\begin{frame}{Why 0.37 / 0.0925 and not 0.4 / 0.1}
  \begin{itemize}\small
  \item Requirement: the pitch must be a whole number of cells on \emph{both} grids.
  \item $3.7 = 37\times0.1$ with 37 prime, so against a 0.1\,mm fine grid the only coarse
        spacings dividing 3.7 exactly are 0.1 and 3.7.
  \item Solving \ti{pitch = n*fine}, \ti{coarse = 4*fine}, $4\mid n$ gives $n=40$:
        fine $=3.7/40=0.0925$, coarse $=3.7/10=0.37$; the validated $4{:}1$ ratio is kept.
  \item Consequence: the pad/gap split quantises identically on both grids ---
        $2.2\,\text{mm}\rightarrow24$ fine / 6 coarse cells (both 2.22\,mm),
        $1.5\,\text{mm}\rightarrow16/4$ (both 1.48\,mm), $24+16=40=4(6+4)$.
  \item \ti{hybrid\_iterate.\_cells()} raises on a fractional cell count, so this is a hard
        constraint of the code, not a convention.
  \end{itemize}
\end{frame}

\begin{frame}[allowframebreaks]{\texttt{run-for-larpix-v2a-wogrid.sh} / \texttt{run-for-largepix-wgrid.sh}}
  \begin{itemize}\small
  \item Same command sequence, explicit \ti{domain} / \ti{gen} / \ti{fdm} instead of
        \ti{field-solve}; one full-depth 0.1\,mm solve per field, no interface.
  \item PART A: \ti{domain --shape=44,44,1601 --spacing '0.1*mm'}; \ti{gen --generator
        pcb\_drift\_pixel\_with\_grid}; \ti{fdm --edges per,per,fix --engine torch
        --nepochs 10 --epoch 130000000 --precision 2e-8 --insulator initial/drift\_insulator}
        $\rightarrow$ \ti{potential/drift3d}.
  \item PART B: identical three commands to \ti{run-pixel-field.sh} PART B.
  \item PART C: \ti{domain --shape=220,220,1601}; \ti{gen --generator pcb\_pixel\_with\_grid};
        \ti{fdm --edges fix,fix,fix --multisteps no --insulator initial/weight\_insulator}
        $\rightarrow$ \ti{potential/weight3d}; then \ti{induce-pixel --npixels 2}.
  \item The two scripts differ in exactly three lines: default store name and the two
        config filenames.
  \item \textbf{With-grid vs without-grid is not a flag}: it follows from
        \ti{GridHoleShape} in the JSON (\ti{"None"}, \ti{"circular"}, \ti{"square"}),
        which is what the generators branch on.
  \end{itemize}
\end{frame}

\begin{frame}{Configs used by \texttt{scripts/}}
  \small
  \begin{center}
  \begin{tabular}{lccc}
    \toprule
    Key & larpix v2a & largepix wgrid & task13 fine \\
    \midrule
    \ti{GridHoleShape}      & \ti{None}   & \ti{circular} & \ti{square} \\
    \ti{HoleRadius}         & 1.9         & 1.9           & 1.9 \\
    \ti{PcbWidth}           & 1.6         & 3.2           & 1.6 \\
    \ti{GridPotential}      & $-2000$     & $-2000$       & $-900$ \\
    \ti{CathodePotential}   & $-7500$     & $-9500$       & $-8400$ \\
    \ti{pixelSize}/\ti{pixelGap} & 3.5 / 0.9 & 3.8 / 0.6  & 2.2 / 1.5 \\
    \ti{chamfer\_r}         & 0.7         & 0.4           & 0.5 \\
    \ti{chamferMode}        & \ti{dynamic}& \ti{dynamic}  & \ti{dynamic} \\
    \ti{padThicknessCells}  & 3           & 3             & 3 \\
    \ti{FR4Thickness}       & 0.1         & 0.1           & 0.1 \\
    \ti{enableInsulatorFR4} & \ti{true}   & \ti{true}     & \ti{true} \\
    \ti{Npixels}            & 5           & 5             & 9 \\
    \ti{driftZDepth}        & 159.9       & 159.9         & 159.9 \\
    \ti{nGridPoints}        & 44          & 10            & 10 \\
    \bottomrule
  \end{tabular}
  \end{center}
\end{frame}

\begin{frame}[allowframebreaks]{Config invariants recorded in the scripts}
  \begin{itemize}\small
  \item \ti{enableFR4} is \textbf{inert} on this branch: no reader in \ti{pochoir/}.
        The pad/FR4 laminate layout is gated on \ti{enableInsulatorFR4} alone
        (\ti{gen\_pcb\_drift\_pixel\_with\_grid.py:507}), which is kept \ti{true}.
  \item The four \ti{task13} configs in \ti{scripts/} are hand-maintained copies of
        same-named files in \ti{test/}; no symlink, no generator, no committed equality
        check. Manual check: \ti{orig=../test ; for f in *task13*.json ; do cmp "\$f" "\$orig/\$f" ; done}.
  \item \textbf{Enforcement-free contract}, stated in the header and enforced by the
        flags in PART B:
    \begin{itemize}
    \item \ti{--insulator} is passed to the \emph{field solve} only;
    \item \ti{velo} receives neither \ti{--boundary} nor \ti{--insulator};
    \item \ti{drift} receives no \ti{--insulator}, so paths are never clamped or
          terminated at a surface;
    \item \ti{--interp-order linear} is required; cubic overshoots at the seam /
          pad-plane kink.
    \end{itemize}
  \item The weighting field is used at its native fine spacing throughout; it is never
        coarsened before the Ramo calculation.
  \end{itemize}
\end{frame}

% =====================================================================
\section{New library modules}
% =====================================================================

\begin{frame}[allowframebreaks]{\texttt{pochoir/nearfar.py} (328 lines)}
  \begin{itemize}\small
  \item Engine-agnostic overlapping-Schwarz near/far solve; the FDM solver arrives as
        two injected callables \ti{solve\_near}, \ti{solve\_far}, each
        \ti{solve(iarr, barr) -> ndarray}.
  \item \ti{resample\_plane} / \ti{resample\_volume}: coordinate-based multilinear
        interpolation with query points clamped to the source bounds, so up- and
        down-sampling are the same code path.
  \item \ti{graft\_plane(init, bmask, ..., fix)}: writes a resampled plane in place and
        optionally marks it Dirichlet.
  \item \ti{seed\_near}: upsample the coarse solution onto the near grid, then re-impose
        exact electrode values at boundary cells.
  \item \ti{\_interface\_indices}: resolves \ti{(ci, ci\_in, nt, nt\_in)}; validates that
        the interface lands on a coarse node, that the near top plane equals the interface,
        and that the far Dirichlet plane lies inside the near $z$-range.
  \item \ti{schwarz\_solve(..., axis=2, interface\_z, tol, max\_iters, near\_start, overlap)}
        returns \ti{(near\_pot, far\_pot, n\_iters, delta)}.
  \end{itemize}
\end{frame}

\begin{frame}[allowframebreaks]{\texttt{nearfar.schwarz\_solve}: one sweep}
  \begin{itemize}\small
  \item \textbf{Far solve}: warm-start from the previous far, re-impose electrode Dirichlet
        values, pin the plane \ti{overlap} coarse cells below the interface to the
        resampled near solution, seed (free) the interface plane, solve the whole coarse
        domain.
  \item \textbf{Near solve}: pin the interface plane to the new far solution, seed (free)
        the adjacent plane, solve the fine near domain.
  \item Convergence measure: \ti{delta = max|near\_new - near\_prev|} against \ti{tol}.
  \item Each sweep ends on the near solve, so the returned near is exactly pinned to the
        returned far at the interface ($C^0$ by construction).
  \item \ti{overlap=1} reproduces the original single-cell coupling byte-for-byte; wider
        overlap converges faster and drives the seam toward gradient continuity.
  \item \ti{near\_start} lets the driver hand in an already-solved sweep-0 near field so its
        intermediate store files are preserved.
  \end{itemize}
\end{frame}

\begin{frame}{\texttt{pochoir/hybrid\_iterate.py} (722 lines) --- sequence}
  \begin{enumerate}\small
  \item Coarse full-volume solve at the coarse spacing.
  \item \ti{refine} the coarse potential onto the fine near grid; \ti{near-bc} pins the
        near grid's last $z$ plane to the coarse value there (Dirichlet, never revisited).
  \item Near solve at the fine spacing on $z=0\ldots$interface, no-flux insulator BC active.
  \item One banded Schwarz sweep at \ti{band\_cells=2} coarse cells --- a 3-coarse-node band
        at the interface. The far solve pins the \emph{inner} node Dirichlet and leaves the
        middle and outer nodes seeded but free; the near then re-solves with its outer node
        pinned. Order: coarse $\rightarrow$ near $\rightarrow$ far $\rightarrow$ near.
  \item \ti{stitch-near} upsamples the Schwarz far field onto the full fine grid and
        overwrites the near planes with the Schwarz near solution. That array \emph{is} the
        final field, written to the profile's output key.
  \end{enumerate}
  \begin{itemize}\small
  \item \ti{max\_sweeps=0} skips step 4 and reproduces the one-shot path exactly.
  \item Every step runs an existing click command via \ti{ctx.invoke}, sharing one store.
  \end{itemize}
\end{frame}

\begin{frame}[allowframebreaks]{\texttt{hybrid\_iterate.py} --- field profiles and keys}
  \small
  \begin{center}
  \begin{tabular}{lll}
    \toprule
    & \ti{drift} & \ti{weighting} \\
    \midrule
    generator & \ti{pcb\_drift\_pixel\_with\_grid} & \ti{pcb\_pixel\_with\_grid} \\
    edges     & \ti{per,per,fix}   & \ti{fix,fix,fix} \\
    transverse& 1 periodic pixel tile & \ti{Npixels} pitches \\
    prefix    & (none)             & \ti{w\_} \\
    \ti{full\_leaf} & \ti{drift3d} & \ti{weight3d} \\
    unit      & V                  & dimensionless \\
    \bottomrule
  \end{tabular}
  \end{center}
  \begin{itemize}\small
  \item Intermediate leaves (\ti{coarse}, \ti{near}, \ti{near\_bc}, \ti{near\_refined},
        \ti{near\_schwarz}, \ti{coarse\_schwarz}) carry the profile prefix.
  \item The full-volume grid is named by field in every taxon:
        \ti{domain/drift3d}, \ti{boundary/drift3d}, \ti{initial/drift3d},
        \ti{potential/drift3d} (likewise \ti{weight3d}).
  \item Namespacing is load-bearing: both fields write into one store so \ti{induce-pixel}
        can read \ti{paths/drift3d} and \ti{potential/weight3d} together, and the
        \ti{\_want} guard must not skip a weighting step because a drift key exists.
  \item No backward compatibility with the retired \ti{fine01}/\ti{w\_fine01} names.
  \end{itemize}
\end{frame}

\begin{frame}[allowframebreaks]{\texttt{hybrid\_iterate.py} --- grids and constants}
  \begin{itemize}\small
  \item \ti{GRID\_SPEC}: \ti{('coarse','coarse','full')}, \ti{('near','fine','near')},
        \ti{(FULL\_LEAF,'fine','full')}; \ti{FULL\_LEAF='fine01'} is a spec-side label only.
  \item \ti{\_extents(prof,cfg,coarse\_spacing,interface\_mm)}: transverse
        $=(\ti{pixelSize}+\ti{pixelGap})\times N_{\text{pix}}$; full depth $=$
        \ti{driftZDepth}$+$\ti{MIN\_CATHODE\_CLEARANCE} rounded up to a whole coarse cell, so
        the cathode does not move with resolution; near depth $=$ interface.
  \item \ti{\_cells(extent,spacing,what,closed)}: transverse is \emph{not} closed
        ($N=$extent/spacing, far node is the wrap of the near one); $z$ \emph{is} closed
        ($N+1$, both faces real planes). Raises on a fractional count.
  \item \ti{\_derive\_grids} (\ti{--domain yes}) vs \ti{\_manual\_grids} (\ti{--domain no},
        all three shapes required).
  \item \ti{\_check\_interface}: the split plane is set by the near grid's $z$ extent, not by
        \ti{--interface}; a mismatch now raises instead of being ignored.
  \item Fixed solver flags: \ti{ENGINE='torch'}, \ti{EPOCH=130000000}, \ti{NEPOCHS=10}.
        Sweep defaults: \ti{DEFAULT\_BAND\_CELLS=2}, \ti{DEFAULT\_MAX\_SWEEPS=1},
        \ti{DEFAULT\_TOL='1*V'}; \ti{DEFAULT\_SPACINGS=\{coarse:0.4, fine:0.1\}}.
  \item \ti{\_want(ctx,targets,thunk)}: the shell resume guard in Python --- skip when all
        targets exist, else run and raise if a target is still missing.
  \end{itemize}
\end{frame}

\begin{frame}[allowframebreaks]{\texttt{pochoir/single\_field.py} (207 lines)}
  \begin{itemize}\small
  \item The same pipeline with the hybrid removed: one grid, one spacing, one solve ---
        \ti{domain/<full\_leaf>} $\rightarrow$ \ti{gen} $\rightarrow$ one \ti{fdm}
        $\rightarrow$ \ti{potential/<full\_leaf>}.
  \item Imports \ti{DEFAULT\_SPACINGS}, \ti{FULL\_LEAF}, \ti{\_cells}, \ti{\_domains},
        \ti{\_extents}, \ti{\_key}, \ti{\_output\_key}, \ti{\_profile}, \ti{\_solve},
        \ti{\_want} from \ti{hybrid\_iterate}; nothing is copied and
        \ti{hybrid\_iterate} is not modified.
  \item Same output keys and same final lattice as the hybrid path, so the two modes are
        comparable plane for plane.
  \item \ti{DEPTH\_QUANTUM = DEFAULT\_SPACINGS['coarse']}: the cathode plane is quantised to
        the coarse spacing even though nothing is solved at that spacing, so the electrode
        sits at a fixed physical depth and the derived grid matches the hybrid fine grid.
  \item \ti{shape} is the primary input and is used verbatim; \ti{derive\_domain=True}
        computes it from the config instead; supplying neither raises.
  \item Only \ti{\_generate\_one} differs from \ti{hybrid\_iterate.\_generate}: it runs
        \ti{gen} for the one full grid rather than looping over three leaves.
  \end{itemize}
\end{frame}

\begin{frame}{\texttt{pochoir/schema.py} (102 lines)}
  \begin{itemize}\small
  \item Wire-Cell field-response object schema, expressed as \ti{namedtuple} classes.
  \item \ti{FieldResponse(planes, axis, origin, tstart, period, speed)}.
  \item \ti{PlaneResponse(paths, planeid, location, pitch)} --- \ti{paths} must be sorted by
        \ti{pitchpos}.
  \item \ti{PathResponse(current, pitchpos, wirepos)}; region
        $=\ti{int(round(pitchpos/pitch))}$, impact $=\ti{pitchpos}-\text{region}\times\ti{pitch}$.
  \item Module docstring records the supported response types (1D, 2D, 2.5D, 3D), the
        Wire-Cell system of units, and the coordinate convention ($x$ counter to nominal
        drift, $y$ up, $z=x\times y$).
  \end{itemize}
\end{frame}

% =====================================================================
\section{Solver core}
% =====================================================================

\begin{frame}{\texttt{fdm\_generic.py}: 2 functions $\rightarrow$ 10}
  \small
  \begin{center}
  \begin{tabular}{ll}
    \toprule
    Function & Role \\
    \midrule
    \ti{edge\_condition} & periodic / fixed halo fill (gains \ti{info\_msg}) \\
    \ti{stencil} & plain Laplace Jacobi (unchanged) \\
    \ti{stencil\_poisson} & Jacobi with RHS, $\nabla^2\phi=-f$ \\
    \ti{stencil\_poisson\_harmonic} & dielectric path, harmonic face permittivities \\
    \ti{neumann\_coeff} & per-cell reciprocal active-neighbour count \\
    \ti{stencil\_poisson\_neumann} & masked no-flux Laplace \\
    \ti{mirror\_masks} & face-centered insulator mirror geometry (superseded) \\
    \ti{mirror\_project} & face-centered mirror application (superseded) \\
    \ti{padplane\_noflux\_geom} & node-centered pad-plane mirror geometry \\
    \ti{padplane\_noflux} & node-centered mirror application \\
    \bottomrule
  \end{tabular}
  \end{center}
\end{frame}

\begin{frame}{Poisson and dielectric stencils}
  \begin{itemize}\small
  \item \ti{stencil\_poisson(array, source, spacing, res)}:
        \[\phi_{\text{new}}=\frac{1}{2N}\sum_{\text{neigh}}\phi-\frac{h^2}{2N}f\]
        \ti{source} is defined on interior cells only; \ti{res} is zeroed on reuse.
  \item \ti{stencil\_poisson\_harmonic(phi, eps, res)}:
        \[\phi=\frac{\sum_d\left(\varepsilon_{i+\frac12}\phi_{i+1}+\varepsilon_{i-\frac12}\phi_{i-1}\right)}
                     {\sum_d\left(\varepsilon_{i+\frac12}+\varepsilon_{i-\frac12}\right)},
          \qquad
          \varepsilon_{i\pm\frac12}=\frac{2\varepsilon_i\varepsilon_{i\pm1}}{\varepsilon_i+\varepsilon_{i\pm1}}\]
  \item Both operate on the whole interior at once and accept a preallocated \ti{res}
        (the solver passes \ti{tmp\_core}), which must be zeroed each call or the
        accumulation diverges.
  \end{itemize}
\end{frame}

\begin{frame}{Masked no-flux stencil}
  \begin{itemize}\small
  \item \ti{neumann\_coeff(solid\_mask, like)}: \ti{coeff} $=1/(\text{non-solid neighbour
        count})$, geometry-only, built once.
  \item \ti{stencil\_poisson\_neumann(phi, solid\_mask, coeff, res)}: each free interior cell
        becomes the average of its non-solid neighbours --- a ghost-cell realisation of
        $\partial\phi/\partial n=0$ on every face of the excluded region. No permittivity
        enters.
  \item Multiplying by the reciprocal (rather than dividing by the count) makes the
        all-\ti{False} mask bit-identical to \ti{stencil\_poisson} with \ti{source=None}:
        \ti{coeff}$\;=1/(2N)=$\ti{norm} uniformly.
  \item \ti{neumann\_coeff} asserts that no free interior cell is fully enclosed by solid
        cells (such a cell is decoupled from the solve); the count is clamped to $\ge1$ so
        solid cells never produce \ti{inf}/\ti{nan}.
  \item All operations are array-level and branch-free, hence \ti{torch.compile}-safe.
  \end{itemize}
\end{frame}

\begin{frame}{Face-centered mirror (\texttt{mirror\_masks} / \texttt{mirror\_project})}
  \begin{itemize}\small
  \item Alternative to the masked stencil: relax the whole volume with plain Laplace, then
        set each insulator \emph{surface} cell to the average of its adjacent LAr cells.
  \item A surface cell is an insulator cell with $\ge1$ LAr (free, non-insulator,
        non-Dirichlet) neighbour. Buried and interior insulator cells are left to plain
        Laplace, so the field is carried below the slab and nothing is shielded.
  \item \ti{mirror\_masks(insulator, lar)} returns \ti{dirs} (per-direction masks),
        \ti{cnt} (LAr-neighbour count), \ti{surf} (\ti{cnt>0}); raises on shape mismatch or
        overlap of the two masks.
  \item \ti{mirror\_project(phi, masks)} is pure and branch-free:
        \ti{where(surf, acc/denom, phi)} with \ti{denom = cnt + (cnt==0)}.
  \item Superseded for the pixel-plane case (see next slide) but retained.
  \end{itemize}
\end{frame}

\begin{frame}[allowframebreaks]{Node-centered pad-plane mirror --- geometry}
  \begin{itemize}\small
  \item \ti{padplane\_noflux\_geom(barr, insulator)} derives the interface from the
        Dirichlet geometry; no hand-crafted mask.
  \item \ti{z\_pad}: the drift-facing face of the contiguous block of \emph{partially}
        Dirichlet planes along the last axis. A fully-Dirichlet plane (cathode) or a
        fully-free plane is not an interface.
  \item \ti{gap2d}: the free (gap) nodes on that plane --- the nodes made no-flux.
  \item \ti{drift\_sign} ($\pm1$): points toward the LAr drift volume, taken as the side
        carrying the fully-Dirichlet cathode plane; if both sides are bounded, the larger
        free gap wins.
  \item Returns \ti{z\_pad}, \ti{gap2d}, \ti{drift\_sign},
        \ti{ghost\_mask} (true only on plane \ti{z\_pad-drift\_sign} at gap nodes),
        \ti{roll\_shift} $=-2\cdot$\ti{drift\_sign}, \ti{axis}.
  \item Shield-grid disambiguation: a grid with apertures is a second, detached partially
        Dirichlet plane, so \ti{partial} becomes non-contiguous. When an \ti{insulator}
        mask is supplied, only the contiguous run carrying the FR4 slab is kept; the grid
        stays an ordinary Dirichlet electrode. With \ti{insulator=None} this is a no-op.
  \item Raises when no partial plane exists, when the partial planes are not contiguous, or
        when the pad plane is too close to the domain edge for a ghost node.
  \end{itemize}
\end{frame}

\begin{frame}{Node-centered pad-plane mirror --- application}
  \begin{itemize}\small
  \item \ti{padplane\_noflux(phi, masks)}: on gap nodes,
        \[\phi[\,\cdot\,,\,z_{\text{pad}}-s\,]\;=\;\phi[\,\cdot\,,\,z_{\text{pad}}+s\,],
          \qquad s=\ti{drift\_sign}\]
        implemented as \ti{where(ghost\_mask, roll(phi, roll\_shift, axis), phi)}.
  \item Effect: the centered derivative
        $-\left(\phi_{z+1}-\phi_{z-1}\right)/2\Delta z$ is exactly zero \emph{at} the pad
        node --- the node the drift integrator samples.
  \item The overwritten node is on the non-drift side, so the drift-side field is never
        corrupted; the transverse $\phi$ profile, hence transverse $E$, is untouched.
  \item Documented distinction from the face-centered mirror: equating an insulator body
        cell to its LAr neighbour zeroes the normal field only at the virtual half-node
        \emph{between} two grid nodes, which the nodal central-difference $E$ never samples.
  \item Pure function (returns a new array), so \ti{torch.compile}-safe;
        \ti{masks['ghost\_mask']} is padded and moved to the device once at set-up.
  \end{itemize}
\end{frame}

\begin{frame}[allowframebreaks]{\texttt{fdm\_torch.py} --- signature and step}
  \begin{itemize}\small
  \item \ti{solve(iarr, barr, periodic, prec, epoch, nepochs, info\_msg=None,
        \_dtype=torch.float64, phi0=None, ctx=None, potential=None, increment=None,
        params=None, epsilon=None, insulator=None)}.
  \item \ti{@torch.compile}d \ti{\_compiled\_step(iarr\_pad, tmp\_core, bi\_core,
        mutable\_core, core, periodic, spacing, source, epsilon, masks)} with a three-way
        branch on \ti{None}-ness (static per run, so one specialisation):
    \begin{itemize}
    \item \ti{masks} present $\rightarrow$ plain Laplace over the whole volume, then
          \ti{padplane\_noflux}; no epsilon;
    \item \ti{epsilon} present $\rightarrow$ \ti{stencil\_poisson\_harmonic};
    \item otherwise $\rightarrow$ \ti{stencil\_poisson}.
    \end{itemize}
  \item Update: \ti{iarr\_pad[core] = bi\_core + mutable\_core*tmp\_core}, then the mirror,
        then \ti{edge\_condition}.
  \item Explicit device (\ti{cuda:0} when available) and \ti{\_dtype} for every tensor;
        \ti{torch.cuda.synchronize()} plus per-epoch and total wall-clock timing.
  \item Returns \ti{(potential, increment)} as \ti{.cpu()} tensors.
  \end{itemize}
\end{frame}

\begin{frame}{\texttt{fdm\_torch.py} --- insulator handling and convergence}
  \begin{itemize}\small
  \item \ti{insulator} is an \textbf{enable signal}; its contents are vestigial. When
        supplied, \ti{padplane\_noflux\_geom(barr, insulator)} re-derives the geometry, the
        ghost mask is padded with \ti{False} and moved to the device, and \ti{epsilon} is
        dropped with a logged notice.
  \item Logged at set-up: \ti{z\_pad}, \ti{drift\_sign}, gap-node count; plus shapes,
        devices and dtypes of \ti{bi\_core}, \ti{mutable\_core}, \ti{tmp\_core},
        \ti{iarr\_pad}.
  \item Convergence check moved from once per epoch to every 1000 steps: \ti{prev} is
        cloned at \ti{istep\%1000==0}, and \ti{maxerr} is computed in float64.
  \item Returns early on \ti{maxerr < prec} and also on \ti{maxerr == 0.0}.
  \item With \ti{insulator=None} and \ti{epsilon=None} the solve is plain Laplace, as
        before.
  \end{itemize}
\end{frame}

\begin{frame}{\texttt{fdm\_torch.py} --- two-step mixed precision}
  \begin{itemize}\small
  \item Selected by \ti{fdm --multisteps yes}; driven from \ti{\_\_main\_\_.fdm}.
  \item Step 1: solve $\nabla^2\phi_0=0$ with the given boundary conditions in
        \ti{torch.float32}; stored as \ti{<potential>\_float32} / \ti{<increment>\_float32}.
  \item Step 2: cast $\phi_0$ to float64 and solve
        $\nabla^2\delta=-\nabla^2\phi_0$ in float64 with \ti{iarr*0} as the initial array;
        stored as \ti{<potential>\_float64\_delta} / \ti{<increment>\_float64\_delta}.
  \item Source construction inside \ti{solve}: pad $\phi_0$, apply \ti{edge\_condition},
        take \ti{stencil(phi0\_)}, then
        \ti{source = -(6/spacing**2)*(s - non\_padded\_phi0)*mutable\_core}; the
        \ti{mutable\_core} factor suppresses the source at fixed boundary cells.
  \item Final: \ti{arr = phi\_0 + delta\_phi}; \ti{err = sqrt(err\_phi0**2 +
        err\_delta\_phi0**2)}.
  \item Debug side-saves: \ti{padded\_phi0}, \ti{stencil}, \ti{source}, \ti{source\_FP32},
        \ti{phi0\_withBC}.
  \end{itemize}
\end{frame}

% =====================================================================
\section{Drift integration}
% =====================================================================

\begin{frame}[allowframebreaks]{\texttt{drift\_numpy.py} --- potential-based drift}
  \begin{itemize}\small
  \item New class \ti{PotentialField(domain, potential, temperature, method, verbose)} and
        entry point \ti{solve\_potential(domain, start, potential, temperature, times,
        method, verbose)}.
  \item The \emph{scalar} potential is interpolated once, and $E$ is obtained by
        differentiating the interpolant, so $E=\nabla\phi$ is curl-free by construction ---
        unlike a component-wise interpolation of a vector $E$ field.
  \item \ti{efield()}: central differences at $h=\tfrac12\Delta x_d$, scaled by
        \ti{units.V}; one-sided (clamped) differences on non-periodic axes.
  \item \ti{v = mu(|E|,T) * E} using \ti{lar.mobility}, the same arithmetic as the
        \ti{velo} command.
  \item Integration: \ti{solve\_ivp(..., method='Radau', rtol=1e-10, atol=1e-10,
        t\_eval=times)}.
  \item Outside the solved domain the velocity is zero --- a data-availability limit, not an
        imposed condition.
  \item \ti{method} follows \ti{drift --interp-order} (\ti{linear} or \ti{cubic});
        \ti{RGI(bounds\_error=False, fill\_value=None)} extrapolates rather than injecting a
        spurious 0.
  \end{itemize}
\end{frame}

\begin{frame}{\texttt{drift\_numpy.py} --- periodic transverse wrap}
  \begin{itemize}\small
  \item Applied in both \ti{Simple} (velocity-field interpolation) and \ti{PotentialField}.
  \item \ti{period[d] = shape[d]*spacing[d]}; \ti{periodic = (0,1)}, the transverse axes of
        the \ti{per,per,fix} tile.
  \item The 0-slice is appended as an extra wrap node on each periodic axis, for the
        potential and for every velocity component.
  \item Positions are taken modulo the period before every field lookup; \ti{inside()}
        skips the periodic axes.
  \item The central-difference samples also wrap, so the gradient at the pad centre (which
        sits on the seam) uses the true neighbour rather than a clamped one-sided value.
  \item Without the wrap, a strong transverse field pushes near-seam electrons past the
        domain edge where \ti{fill\_value=0} freezes them, instead of letting them wrap onto
        the pad centre.
  \end{itemize}
\end{frame}

\begin{frame}{\texttt{drift\_numpy.py} --- removed enforcement}
  \begin{itemize}\small
  \item The module docstring records that all drift-side enforcement is deleted:
    \begin{itemize}
    \item the terminal event that force-stopped a path at the FR4 top face;
    \item the $v=0$ zeroing inside the insulator;
    \item the mask-aware (ghost/mirror) correction of $E$ at the FR4 surface.
    \end{itemize}
  \item The drift now integrates $v=\mu\,\nabla\phi$ on the solved potential; $E$ is taken
        plainly, with no insulator mask or field correction at this stage.
  \item The only insulator condition anywhere is the no-flux Neumann boundary the Laplace
        solver imposed when producing the potential.
  \item \ti{DRIFT\_NONE = 0} is retained and always returned as \ti{endtag}, for call-site
        compatibility.
  \item Hardcoded-index diagnostic prints in \ti{solve\_sde} are commented out.
  \end{itemize}
\end{frame}

% =====================================================================
\section{CLI}
% =====================================================================

\begin{frame}{\texttt{\_\_main\_\_.py}: 1545 $\rightarrow$ 2658 lines}
  \begin{itemize}\small
  \item Module-level \ti{logging} to \ti{\$POCHOIR\_STORE/pochoir.log}
        (\ti{POCHOIR\_LOG} overrides the path), exposing
        \ti{info\_msg} / \ti{err\_msg} / \ti{debug\_msg}.
  \item \ti{STORE\_DIR} honours \ti{POCHOIR\_STORE}, so side-saved PNGs and legacy
        \ti{.npy} dumps follow the renamed output folder instead of a hardcoded
        \ti{store/}; the directory is created at import.
  \item \ti{pochoir.hybrid\_iterate} and \ti{pochoir.single\_field} are imported at module
        scope so click option \emph{defaults} come from the module constants rather than
        repeated literals.
  \item New helpers: \ti{\_coarse\_interpolator}, \ti{\_interp\_clamped},
        \ti{make\_pixel\_start\_points}, \ti{\_shift\_paths\_pixel\_grid},
        \ti{\_load\_pixel\_geometry}, \ti{\_load\_start\_point\_config},
        \ti{\_field\_solve\_given}, \ti{\_field\_solve\_opt}.
  \end{itemize}
\end{frame}

\begin{frame}{New commands: near/far primitives}
  \begin{itemize}\small
  \item \ti{refine -c COARSE -i INITIAL -b BOUNDARY -I OUT}: multilinear upsample of a
        coarse potential onto the fine grid using actual grid coordinates, then merge exact
        fine boundary values at boundary cells. Output taxon \ti{initial}.
  \item \ti{near-bc -i INITIAL -b BOUNDARY -c COARSE -I IOUT -B BOUT [--axis 2]}: pin the
        last plane along \ti{--axis} Dirichlet to the interpolated coarse potential; emits a
        new \ti{initial} and \ti{boundary} pair.
  \item \ti{coarsen -i INPUT -d DOMAIN -o OUT [--average]}: downsample onto a coarser
        domain by stride sampling, or block-average over stride-sized blocks; validates the
        result against the target domain shape.
  \item \ti{stitch-near -n NEAR -c COARSE -d DOMAIN -I OUT [--axis 2]}: upsample the coarse
        field onto the full fine domain, overwrite the first \ti{near.shape[axis]} planes
        with the near solution, and stamp \ti{domain} into the output metadata so \ti{velo}
        resolves the grid downstream.
  \end{itemize}
\end{frame}

\begin{frame}{New command: \texttt{near-far-solve}}
  \begin{itemize}\small
  \item Inputs: \ti{--coarse-potential}, \ti{--coarse-initial}, \ti{--coarse-boundary},
        \ti{--near-initial}, \ti{--near-boundary}, \ti{--near-potential} (warm start),
        \ti{--insulator}.
  \item Geometry: \ti{--interface}, \ti{--axis} (default 2), \ti{--edges}.
  \item Solver: \ti{--engine}, \ti{--epoch} (default $10^6$), \ti{--nepochs} (10),
        \ti{--near-precision} (\ti{2e-11}), \ti{--far-precision} (\ti{2e-7}).
  \item Schwarz: \ti{--tol} (units accepted, e.g. \ti{'1*V'}), \ti{--max-iters} (6),
        \ti{--overlap} (1, in coarse cells).
  \item Outputs: \ti{--near-out}, \ti{--far-out} (both required).
  \item \ti{--insulator} is applied to the \emph{near} re-solves only; the coarse/far
        re-solves never see it, as FR4 is unresolved at the coarse pitch.
  \item Delegates to \ti{nearfar.schwarz\_solve} with \ti{solve\_near}/\ti{solve\_far} bound
        to \ti{pochoir.fdm.solve\_<engine>}.
  \end{itemize}
\end{frame}

\begin{frame}[allowframebreaks]{New commands: \texttt{hybrid-iterate} and \texttt{field-solve}}
  \begin{itemize}\small
  \item \ti{hybrid-iterate}: \ti{--coarse-config}, \ti{--fine-config}, \ti{--interface},
        \ti{--precision}, \ti{--field \{drift,weighting\}}, \ti{--coarse-spacing 0.4},
        \ti{--fine-spacing 0.1}, \ti{--domain \{yes,no\}} (default \ti{yes}),
        \ti{--coarse-shape}, \ti{--near-shape}, \ti{--fine-shape}, \ti{--band-cells},
        \ti{--max-sweeps}, \ti{--schwarz-tol}.
  \item \ti{field-solve}: single entry point. \ti{--hybrid yes|no} is the only mode switch
        and selects which driver runs, hence which options are legal.
    \begin{itemize}
    \item hybrid-only: \ti{--coarse-config}, \ti{--fine-config}, \ti{--interface}
          (default \ti{'40*mm'}), \ti{--coarse-spacing}, \ti{--fine-spacing},
          \ti{--coarse/near/fine-shape}, \ti{--band-cells}, \ti{--max-sweeps},
          \ti{--schwarz-tol};
    \item single-only: \ti{--config}, \ti{--spacing}, \ti{--shape}.
    \end{itemize}
  \item \ti{--domain} defaults to \ti{no} here (shape options required), unlike
        \ti{hybrid-iterate}'s own \ti{yes}; with \ti{--domain yes} the shape options are
        rejected as contradictory.
  \item Options belonging to the other mode are \emph{rejected}, not ignored: which options
        the user actually typed is read from click's
        \ti{ParameterSource.COMMANDLINE}.
  \end{itemize}
\end{frame}

\begin{frame}[allowframebreaks]{Changed commands}
  \begin{itemize}\small
  \item \ti{fdm}: adds \ti{--epsilon} (permittivity array), \ti{--insulator} (no-flux mask,
        no epsilon), \ti{--multisteps Y/N}; engine choice
        \ti{numpy|numba|torch|cupy|cumba}. The insulator kwarg is passed only when present,
        so other engines and flag-off runs are unchanged.
  \item \ti{velo}: adds \ti{--boundary} and \ti{--insulator}, both documented as accepted
        for compatibility and \textbf{inert} --- velocity stays pure $\mu\,\nabla\phi$.
  \item \ti{drift}: adds \ti{--insulator} (also inert), \ti{--plot}, and
        \ti{--interp-order \{linear,cubic\}} (default \ti{linear}). Resolves the potential
        and temperature from the velocity array's metadata; stores lacking that metadata
        fall back to legacy component-wise interpolation.
  \item \ti{starts}: adds \ti{-c/--config} (multiple) and \ti{--plot}; \ti{-m yes} now
        builds a cell-centred $n\times n$ grid via \ti{make\_pixel\_start\_points} instead of
        a hardcoded array.
  \item \ti{induce-pixel}: adds \ti{-c/--config} (multiple) and \ti{--plot}; per-pad
        re-centring via \ti{\_shift\_paths\_pixel\_grid}; the \ti{sys.exit()} on the induce
        path is removed.
  \end{itemize}
\end{frame}

\begin{frame}{Config-driven start points}
  \begin{itemize}\small
  \item \ti{make\_pixel\_start\_points(z\_depth, ngridpoints, pitch, spacing)} returns
        exactly $n^2$ points at one fixed $z$.
  \item \ti{spacing} defaults to \ti{pitch/ngridpoints}; the first point sits at
        \ti{spacing/2}, so the grid is cell-centred
        (e.g.\ $4.4/10\rightarrow$ 0.22, 0.66, 1.10, \ldots, 4.18).
  \item \ti{\_load\_start\_point\_config} merges all \ti{-c} JSON files and maps:
        \ti{driftZDepth}$\rightarrow$\ti{z\_depth},
        \ti{nGridPoints}$\rightarrow$\ti{ngridpoints},
        \ti{gridSpacing}$\rightarrow$\ti{spacing} (optional),
        \ti{pixelSize}$+$\ti{pixelGap}$\rightarrow$\ti{pitch}.
  \item The same config files are handed to \ti{gen}, \ti{starts} and \ti{induce-pixel}.
  \end{itemize}
\end{frame}

% =====================================================================
\section{Geometry generators}
% =====================================================================

\begin{frame}{\texttt{gen\_pcb\_*\_with\_grid.py} --- shield plane}
  \begin{itemize}\small
  \item \ti{draw\_pcb\_plane(shape, arr, barr, z, r1, gridPotential)}: rewritten
        vectorised --- set the whole plane to boundary at \ti{gridPotential}, then clear
        four quarter-disk holes at the corners with an \ti{mgrid} mask. Writes both
        \ti{arr} and \ti{barr}.
  \item \ti{draw\_pcb\_plane\_rounded\_sq\_drift(arr, barr, p\_gap, p\_size, pcb\_width,
        pp\_loweredge, gridPotential, chamfer\_r)}: solid plane, four rectangular quadrants
        cleared to open square apertures, corners restored by
        \ti{\_apply\_rounded\_corners(..., val=1)}.
  \item Selected by \ti{GridHoleShape}: \ti{"None"} (no shield grid), \ti{"circular"},
        \ti{"square"}.
  \item The shield is a single $z$ plane at \ti{pp\_loweredge + pcb\_width}.
  \end{itemize}
\end{frame}

\begin{frame}[allowframebreaks]{\texttt{gen\_pcb\_*\_with\_grid.py} --- pads and chamfer}
  \begin{itemize}\small
  \item \ti{trimCorner(arr, x, y, z1, z2, corner, val=0, chamfer\_r=4)}: parameterised
        quarter-disk notch at one inner corner; \ti{val=0} carves a pad, \ti{val=1} fills a
        void; \ti{chamfer\_r}$\le0$ is a no-op. Replaces the duplicated
        \ti{trimCorner}/\ti{trimCorner\_pcb} pair.
  \item \ti{trimCorner\_forpix}: the original fixed-cell staircase notch, kept
        byte-identical to the \ti{for\_pix} branch and used only when
        \ti{chamferMode == "fixed\_cell"}. Its reach is 4 grid cells, so its physical size
        scales with the spacing.
  \item \ti{\_apply\_rounded\_corners(barr, p\_size, p\_gap, z1, z2, val, chamfer\_r)}:
        encodes the four corner positions and quadrant indices once, shared by pad and
        shield.
  \item \ti{half = (p\_size+1)//2} replaces \ti{int(p\_size/2)}, fixing the broken
        quarter-pixel corner symmetry.
  \item \ti{draw\_pixel\_plane(..., chamfer\_r=0.7, chamferMode='dynamic', fr4\_bottom=False,
        n\_fr4=0)}: fills the four pad quadrants over the pixel-plane $z$ range, then carves
        the corners with \ti{val=0}.
  \end{itemize}
\end{frame}

\begin{frame}[allowframebreaks]{FR4 laminate and pad thickness}
  \begin{itemize}\small
  \item \ti{fr4\_bottom} / \ti{n\_fr4}: when the pixel-plane laminate is active, the pad
        conductor covers only the \emph{top} of the \ti{pixelPlaneWidth} layer; the bottom
        \ti{n\_fr4} cells stay free so their FR4 permittivity enters the solve. With
        \ti{fr4\_bottom=False} (\ti{n\_fr4=0}) the layout is byte-identical to the previous
        behaviour.
  \item The laminate layout is gated on \ti{enableInsulatorFR4} alone
        (\ti{gen\_pcb\_drift\_pixel\_with\_grid.py:507}); \ti{enableFR4} has no reader in
        \ti{pochoir/}.
  \item \ti{padThickness} / \ti{padThicknessCells}: give the pad real thickness (3 cells in
        the shipped configs) rather than a one-cell sheet.
  \item \ti{gen} also writes the no-flux mask under \ti{<initial>\_insulator}; the solver's
        \ti{--insulator} flag is an enable signal and the interface itself is re-derived by
        \ti{padplane\_noflux\_geom}.
  \item Diagnostic 3-D scatter PNGs of \ti{arr} and \ti{barr} (full and clipped to the first
        150 $z$ planes, at three marker sizes) are written to \ti{\_plot\_dir()}, which
        follows \ti{POCHOIR\_STORE}.
  \end{itemize}
\end{frame}

\begin{frame}[allowframebreaks]{Config key surface}
  \small
  \begin{center}
  \begin{tabular}{ll}
    \toprule
    Key & Meaning \\
    \midrule
    \ti{GridHoleShape} & \ti{None} / \ti{circular} / \ti{square}: selects the shield drawing \\
    \ti{chamfer\_r} & corner radius, mm (quantised to cells at the run's spacing) \\
    \ti{chamferMode} & \ti{dynamic} (mm-based) or \ti{fixed\_cell} (\ti{for\_pix} replica) \\
    \ti{padThickness}, \ti{padThicknessCells} & pad conductor thickness \\
    \ti{FR4Thickness} & laminate thickness under the pad \\
    \ti{enableInsulatorFR4} & gates the pad/FR4 laminate $z$ layout \\
    \ti{enableFR4} & present in configs, no reader in \ti{pochoir/} \\
    \ti{driftZDepth} & electron launch node; also sets the derived cathode plane \\
    \ti{nGridPoints} & start-point grid per side \\
    \ti{LArPermittivity}, \ti{FR4Permittivity} & dielectric path; dropped from recent configs \\
    \bottomrule
  \end{tabular}
  \end{center}
\end{frame}

% =====================================================================
\section{Small fixes and supporting material}
% =====================================================================

\begin{frame}{Compatibility fixes}
  \begin{itemize}\small
  \item \ti{arrays.gradient}: numpy $\ge1.23$ no longer accepts a single 1-D array as
        ``one scalar per axis''. A length-$N$ sequence is now unpacked into $N$ positional
        scalars; the torch branch, which passed the list unsplatted, is fixed the same way.
        The debug print is removed.
  \item \ti{plots.quiver}: \ti{fig.gca(projection='3d')} $\rightarrow$
        \ti{fig.add\_subplot(projection='3d')} for matplotlib $\ge3.5$.
  \item \ti{main.Main.put}: prints the resolved key.
  \item \ti{fdm\_generic.edge\_condition} gains an \ti{info\_msg} keyword.
  \item \ti{.gitignore}: adds \ti{*PLAN*}, \ti{env/}, \ti{toy/*.npy} and the
        \ti{test/} scratch directories (\ti{OUTPUT/}, \ti{bkup\_store/}, \ti{plots/},
        \ti{toy/}, \ti{review/}, log files).
  \end{itemize}
\end{frame}

\begin{frame}{\texttt{README\_for\_pixelatedReadout.md} (220 lines)}
  \begin{itemize}\small
  \item Companion to \ti{README.org}, scoped to \ti{scripts/run-pixel-field.sh}.
  \item Numbers the six pipeline stages and the store key each writes; records that
        \ti{date} is printed at three checkpoints.
  \item Documents the two modes, that only PARTS A and C differ, and that both write the
        same keys and land on the same final lattice.
  \item States the enforcement-free contract and the \ti{--interp-order linear}
        requirement.
  \item Requirements: Linux, Python 3.9+, CUDA GPU strongly recommended,
        \ti{pip install -e .[torch,plots,hdf5]}.
  \item Config table mapping the four shell variables to their files and roles, plus the
        duplicated-config warning and the \ti{cmp} recipe.
  \end{itemize}
\end{frame}

\begin{frame}[allowframebreaks]{\texttt{docs/examinations/wogrid/} (6 documents)}
  \begin{itemize}\small
  \item \ti{00-overview.md} --- what the branch is; changes vs master across the solver,
        stencils, generators, units, LAr physics, tooling.
  \item \ti{01-algorithm.md} --- Jacobi refresher; \ti{stencil\_poisson};
        the two-step mixed-precision solver step by step; convergence cadence; pixel/PCB
        drawing algorithm; domain sizes; diffusion physics; unit-system change.
  \item \ti{02-bugs.md} --- catalogue with severities: docstring/code sign mismatch;
        hardcoded \ti{6} assuming 3-D; \ti{err=None} crash when \ti{epoch<1000};
        \ti{barr\_pad} and \ti{iarr\_pad\_source} allocated but never read;
        inconsistent \ti{mutable\_core} on the source; \ti{numpy.bool} removal.
  \item \ti{03-gpu-efficiency.md} --- \ti{@torch.compile} disabled;
        \ti{torch.cuda.synchronize()} sites; \ti{prev} cloning; per-step kernel-launch
        count; mixed-precision profile; dead profiler code; geometry-generation cost.
  \item \ti{04-memory.md} --- per-tensor inventory and totals for standard and two-step
        modes; memory held outside the solver; reduction opportunities.
  \item \ti{05-summary.md} --- prioritised bug list and recommendations.
  \end{itemize}
\end{frame}

\begin{frame}[allowframebreaks]{Test suite: 29 new \texttt{test/test\_*.py}}
  \begin{itemize}\small
  \item Schwarz and near/far: \ti{test\_hybrid\_schwarz\_step},
        \ti{test\_hybrid\_schwarz\_wiring}, \ti{test\_hybrid\_cli\_schwarz\_options},
        \ti{test\_nearfarsolve}, \ti{test\_near\_far\_solve\_tol},
        \ti{test\_nearfar\_ckpt\_key}.
  \item Solver and boundary: \ti{test\_stencil\_poisson},
        \ti{test\_insul\_bc\_torch\_solve}, \ti{test\_neumann\_bc\_doc}.
  \item Geometry and config: \ti{test\_grid\_init\_method},
        \ti{test\_pixel\_geometry\_loaded\_from\_config}, \ti{test\_start\_point\_config},
        \ti{test\_make\_pixel\_start\_points}, \ti{test\_shift\_paths\_pixel\_grid}.
  \item Task drivers: \ti{test\_task13\_oneshot\_hybrid},
        \ti{test\_task13\_phase1\_geometry\_configs},
        \ti{test\_task13\_runner\_band\_flags}, \ti{test\_task13\_config\_descriptions},
        \ti{test\_task1\_fixedcell\_2cm\_driver}, \ti{test\_task9\_drift\_schwarz},
        \ti{test\_run\_full\_3d\_pixel\_driver}, \ti{test\_drift\_paths\_task10\_notebook}.
  \item Hygiene: \ti{test\_no\_stray\_prints}, \ti{test\_pipeline\_no\_todo\_markers},
        \ti{test\_notebooks\_have\_no\_outputs},
        \ti{test\_notebook\_enforcementfree\_halt},
        \ti{test\_induce\_pixel\_no\_sys\_exit}, \ti{test\_plot\_flag},
        \ti{test\_hybrid\_docs\_describe\_sweep}.
  \end{itemize}
\end{frame}

\begin{frame}[allowframebreaks]{Other supporting material}
  \begin{itemize}\small
  \item \ti{Makefile}: \ti{make test-pixel} runs pytest and the pipeline, capturing
        \ti{pytest.log}, \ti{pipeline.log}, \ti{outputs/} and a \ti{NOTES.md} stub under
        \ti{test/review/<date>\_<slug>/}. Shapes default to \ti{44,44,300} /
        \ti{220,220,300}, overridable via \ti{POCHOIR\_DRIFT\_SHAPE} /
        \ti{POCHOIR\_WEIGHT\_SHAPE}.
  \item \ti{test/} runners: \ti{run-task1} through \ti{run-task13}, plus
        \ti{run-full-3d-pixel.sh}, \ti{run-full-3d-pixel-debug.sh},
        \ti{run-fulldepth-01mm-neumann.sh}, \ti{run-hybrid-\{10,30\}cm-drift.sh},
        \ti{run-near2cm-drift.sh} and the \ti{test-full-3d-*} variants.
  \item Analysis helpers: \ti{bench\_fdm.py}, \ti{parse\_maxerr.py},
        \ti{visualize\_fdm\_domain.py}, \ti{sample\_field\_along\_paths.py},
        \ti{analyze\_gap\_endpoint\_\{efield,velocity\}.py}.
  \item Notebooks: \ti{test/DriftPaths\_task10\_2cm\_01mm.ipynb},
        \ti{PathPixBorder\_neumannFullDetph\_01mm.ipynb},
        \ti{Task10\_surface\_fields\_investigation.ipynb},
        \ti{test/for\_pixel/*.ipynb}, \ti{toy/*.ipynb}.
  \item \ti{toy/InverseDistanceWeight.py} and \ti{InverseDistanceWeight\_torch.py}.
  \item Write-ups: \ti{test/neumann\_bc\_fr4\_insulator.tex},
        \ti{test/code\_structure\_test\_full\_3d\_pixel.md},
        \ti{test/run-full-3d-pixel.md}.
  \end{itemize}
\end{frame}

% =====================================================================
\section{Head commit}
% =====================================================================

\begin{frame}[allowframebreaks]{Head commit \texttt{532e923}: task13 retarget}
  \begin{itemize}\small
  \item Geometry: pitch 3.7\,mm (\ti{pixelSize} 2.2, \ti{pixelGap} 1.5),
        \ti{Npixels} 9, \ti{driftZDepth} 159.9, \ti{CathodePotential} $-8400$,
        \ti{GridHoleShape} \ti{square} at \ti{GridPotential} $-900$, \ti{chamfer\_r} 0.5.
  \item Spacings: coarse 0.37, fine 0.0925 (see the derivation slide); the $4{:}1$ ratio is
        preserved and the pad quantises identically on both grids, so the two configs no
        longer carry different pad/gap splits.
  \item \ti{chamfer\_r} 0.5 quantises to 5 cells $=0.4625$\,mm at 0.0925\,mm --- a real
        0.0375\,mm change to the pad chamfer footprint relative to the retired 0.1\,mm runs.
  \item \ti{run-pixel-field.sh} carries the re-derived shape table and a note that
        \ti{--domain no} skips the \ti{\_cells()} check.
  \item Cathode formula recorded in the configs:
        \ti{CathodePotential = -BulkField * (driftZDepth - pixelPlaneLowEdgePosition)};
        $-8400$\,V over 149.9\,mm is 56.04\,V/mm nominal, maintained by hand.
  \item Corrected claim: the bulk field is set by the cathode against the pixel plane at
        0\,V, not grid-to-cathode --- the shield grid has holes and is transparent at tile
        centre (measured $\approx-71$\,V at the grid plane on the earlier
        149.9\,mm / $-7500$\,V run).
  \item Task13 tests now read the cathode target and its worked result out of the config
        prose and check self-consistency against each file's own geometry; the
        byte-identical-keys check reads the configs at a fixed commit rather than the
        working tree.
  \end{itemize}
\end{frame}

\end{document}
